\documentclass[11pt]{article}
\usepackage[margin=1in]{geometry}
\usepackage{amsmath,amssymb,amsthm,booktabs}
\usepackage[hidelinks]{hyperref}
\setlength{\emergencystretch}{2em}
\newtheorem{theorem}{Theorem}
\newtheorem{proposition}[theorem]{Proposition}
\newtheorem{corollary}[theorem]{Corollary}
\newcommand{\Z}{\mathbb Z}
\newcommand{\F}{\mathbb F}
\newcommand{\C}{\mathbb C}
\newcommand{\rank}{\operatorname{rank}}
\title{Simple-module bounds and exhaustion of the single-$A_5$ cocycle route}
\author{Internal research working proof}
\date{16 November 2026}
\begin{document}
\maketitle

For every retained double-branched-cover group $G$, every field $k$
and every finite-dimensional module $V$ whose action factors through
a homomorphism $G\to A_5$, the new complete coefficient audit gives
\begin{equation}\label{eq:bound}
 \dim_k Z^1(G,V)\le\frac32\dim_k V.
\end{equation}
Thus the unrestricted linear-cocycle estimate cannot force
$\rank G\ge3$ through any single $A_5$ quotient on the sixteen-target
panel, even with nonsemisimple coefficients of arbitrarily large
dimension. This is a stronger boundary than the earlier tests of
the two permutation augmentation modules. It does not prove that
$G$ is two-generated or settle a target knot.

The argument combines an elementary composition-series reduction
with the remaining three-dimensional and characteristic-two
two-dimensional simple modules. All sixteen quotient-kernel classes
on twelve knots are checked, with the four no-map targets handled
separately. The absolute simple-module classification is imported
from ATLAS \cite{atlas}; the supplied models and finite matrices have
exact independent replay. Uniform proofs are unformalized and
externally unreviewed, with no priority claim or full MRC conclusion.

\section{Why ordinary coefficient searches reduce to simple modules}

Let $G$ be finitely presented, $\phi:G\to Q$ a fixed finite quotient
and $V$ a finite-dimensional $kQ$-module, regarded as a $G$-module.
If $G$ has $d$ generators, evaluation on them injects $Z^1(G,V)$
into $V^d$. Hence
\[
 \rank G\ge\left\lceil\frac{\dim Z^1(G,V)}{\dim V}\right\rceil,
 \qquad
 \dim Z^1(G,V)=\dim H^1(G,V)+\dim V-\dim V^G.
\]

\begin{proposition}\label{prop:simple}
For a fixed field and quotient action, the largest ratio
$\dim Z^1(G,V)/\dim V$ over nonzero finite-dimensional modules is
attained on a simple $kQ$-module.
\end{proposition}
\begin{proof}
For $0\to U\to V\to W\to0$, the induced map
$Z^1(G,V)\to Z^1(G,W)$ has kernel $Z^1(G,U)$. Consequently
\[
 \dim Z^1(G,V)\le\dim Z^1(G,U)+\dim Z^1(G,W).
\]
Apply this repeatedly along a composition series. If its factors
are $S_i$, the ratio is bounded by the weighted average of
$\dim Z^1(G,S_i)/\dim S_i$, hence by their maximum. Simple modules
themselves are allowed, giving equality of the maxima.
\end{proof}

This includes nonsplit extensions; semisimplicity is unnecessary.
The bound may be checked after extending scalars to an algebraic
closure: the finite relator matrix defining $Z^1$ commutes with
field extension, preserving its dimension and that of $V$.
The proposition concerns unrestricted ordinary-group cocycles.
It does not replace meridional estimates whose evaluation spaces
are smaller than $V$, or exclude nonlinear generation obstructions.

\section{The remaining simple modules and exact calculations}

The absolute simple degrees needed here are
\[
\begin{array}{c|l}
 \operatorname{char}k&\text{degrees}\\\hline
 0\text{ or }p\nmid60&1,3,3,4,5\\
 2&1,2,2,4\\
 3&1,3,3,4\\
 5&1,3,5.
\end{array}
\]
The modular lists are explicitly supplied by ATLAS \cite{atlas}.
All nontrivial simples are faithful since $A_5$ is simple.
The cross-characteristic list is the ordinary character list under
the standard semisimple reduction away from $60$. Completeness is
an imported representation-theoretic input, not an inference from
a bounded module search.

Write $\mathcal O=\Z[\tau]$, $\tau^2=\tau+1$. Rotation matrices
from the retained integral $H_3$ model give the two three-dimensional
simple modules through the two embeddings of $\mathcal O$.
For every retained map use the thirteen-generator normalized
core presentation, including its fourteen relators. Full cocycles
have a $42$ by $39$ matrix over $\mathcal O$. No knot-group extension,
meridian-root condition or deck-involution constraint is imposed.

Expanding multiplication by $a+b\tau$ as
$\left(\begin{smallmatrix}a&b\\b&a+b\end{smallmatrix}\right)$ gives
an $84$ by $78$ integer matrix. Its nonzero Smith factors are
\[
 \underbrace{1,\ldots,1}_{71},5
 \quad\text{on eight classes},\qquad
 \underbrace{1,\ldots,1}_{70},7,35
 \quad\text{on eight classes}.
\]
The expanded invariant-vector matrix has factors $1,1,1,1,2,2$.
Thus all possible rank changes are isolated exactly. The ramified
prime $5$ is evaluated separately at $\tau=3$, rather than dividing
a regular-expansion rank by two. Primes $2,3,7$ are inert and are
tested over the corresponding quadratic residue fields.

The resulting three-dimensional data are
\[
\begin{array}{c|c|c|c}
 \operatorname{char}k&\dim Z^1&\dim V^G&\dim H^1\\\hline
 0\text{ or }p\notin\{2,5,7\}&3&0&0\\
 2&3&1&1\\
 5&4&0&1\\
 7\text{ on the eight }(7,35)\text{ classes}&4&0&1\\
 7\text{ on the other eight classes}&3&0&0.
\end{array}
\]
In characteristic two the three-dimensional model has an invariant
line. Its quotient and the quotient's Frobenius conjugate are the
two absolute two-dimensional simples over $\F_4$. Both have
$\dim Z^1=3$, no invariants and $\dim H^1=1$ on every retained map.
The invariant line has no invariant complement, checked by the
absence of fixed covectors; this also provides an explicit nonsplit
module in the scope of Proposition~\ref{prop:simple}.

The earlier five- and six-point cover receipts supply $V_4$ and
$V_5$. For $V_4$ in characteristic other than five, the cocycle
dimension is four, except that it is five at three. For $V_5$ in
characteristic other than two and three, it is five, except that it
is six at seven on the same eight exceptional classes. In particular
the required simple $V_4$ at two and three, and simple $V_5$ at five,
are covered. Full matrix-algebra spans verify absolute simplicity
of these modular models. The trivial module has cocycle dimension
at most one, since each $G^{\rm ab}=\Z/D$.

The largest ratios on absolute simples are consequently
\[
\begin{array}{c|c}
 \operatorname{char}k&\text{largest ratio}\\\hline
 2&3/2\\
 3&5/4\\
 5&4/3\\
 7\text{ on the eight exceptional classes}&4/3\\
 \text{all other cases, including characteristic zero}&1.
\end{array}
\]
Proposition~\ref{prop:simple} proves equation~\eqref{eq:bound} for
all modules and fields. The factor $3/2$ is attained on the
two-dimensional characteristic-two simples.

\section{Scope over maps and over induced covers}

The earlier complete kernel enumeration covers every epimorphism
to $A_5$, modulo automorphisms, on this panel. Postcomposition by
an automorphism only permutes the coefficient modules, so the bound
is independent of representative. Every base has cyclic
abelianization of order coprime to $30$. A proper nontrivial subgroup
of $A_5$ has a nontrivial abelian quotient at $2$, $3$ or $5$;
thus a nontrivial homomorphism from $G$ to $A_5$ is already onto.
The four no-map targets therefore admit only the trivial such
action, which also satisfies the bound.

\begin{corollary}\label{cor:covers}
Let $H=\phi^{-1}(L)$ for any subgroup $L\le A_5$, and put
$n=[G:H]$. Over every field,
\[
 b_1(H;k)\le\lfloor n/2\rfloor+1.
\]
Consequently none of these induced covers can cross the Schreier
threshold for proving $\rank G\ge3$. When the characteristic does
not divide $n$, its augmentation-module Betti-gain test also fails.
Moreover $b_1(H;\C)=0$, including the degree-sixty kernel covers.
\end{corollary}
\begin{proof}
For the transitive permutation module $M=k[G/H]$, Shapiro gives
$H^1(G,M)=H^1(H,k)$, while $\dim M^G=1$. Equation~\eqref{eq:bound}
therefore gives $b_1(H;k)\le n/2+1$. To force rank three by
Schreier would require $b_1(H;k)\ge n+2$.

If $n$ is invertible, its $(n-1)$-dimensional augmentation summand
$A$ has no invariants. Thus the Betti gain equals
$\dim H^1(G,A)\le(n-1)/2$, whereas rank three requires gain at
least $n$. Finally, in characteristic zero every simple has zero
$H^1$ by the exact data above and the earlier $V_4,V_5$ homology
data. All modules are semisimple, so $H^1(G,M)=0$.
\end{proof}

This covers every $A_5$-induced degree, not every cover of that
degree with other monodromy. It is an exclusion of these estimates,
not an upper certificate for the actual group rank.

\section{Two quotient kernels and independent verification}

A module for a joint quotient $A_5\times A_5$ need not factor through
either projection. On the four target knots with two retained
kernels, their joint image has order $3600$. The four Frobenius
choices of the natural tensor $2\otimes2$ give sixteen absolutely
irreducible four-dimensional actions over $\F_4$. All have
\[
 \dim Z^1=6,\qquad\dim H^1=2.
\]
They give the ordinary lower bound two; a sharp-four certificate
would require $\dim Z^1\ge9$, equivalently $\dim H^1\ge5$.
Other product modules and characteristics remain untested.

The independent verifier rebuilds the integral $H_3$ representation
as ordinary six-by-six matrices, evaluates affine relator words on
points from right to left, and replays every Smith operation. It
checks each residue dimension, the invariant-line quotient,
Frobenius conjugation, full modular matrix-algebra spans, joint
image orders and all sixteen tensor cocycle matrices. It uses no
producer quadratic-field kernel solver or Fox-prefix constructor.
The PD parser, original input maps, old coefficient-group convention
and Smith replay are shared and explicit. Free-group controls recover
$\dim Z^1(F_d,V)=d\dim V$ for $d=2,3$.

The complete receipt is \path{a5_all_modules_verification.json};
plans and all matrices are retained. Further linear work needs a
different quotient or an untested product-module structure, not
larger coefficient modules through a single retained $A_5$ map.
The restricted meridional geometry and genuinely different
nonlinear rank obstructions remain separate routes. No new sharp
knot case, actual rank-two certificate or full MRC solution is claimed.

\begin{thebibliography}{9}
\bibitem{atlas} ATLAS of Finite Group Representations,
\emph{Alternating group $A_5$, linear groups $L_2(4)$ and $L_2(5)$},
representation lists in characteristics $0,2,3,5$.
\url{https://brauer.maths.qmul.ac.uk/Atlas/alt/A5/}.
Primary snapshot retained with SHA256, accessed 16 November 2026.
\end{thebibliography}
\end{document}
